\documentclass[11pt,a4paper]{article}
\usepackage{fontspec}
\usepackage{babel}
\babelprovide[import,main]{russian}
\babelfont{rm}{Noto Serif}
\babelfont{sf}{Noto Sans}
\babelfont{tt}{Noto Sans Mono}
\usepackage[a4paper,margin=23mm,headheight=14pt,headsep=8mm]{geometry}
\usepackage{amsmath,amssymb,amsthm,mathtools,stmaryrd}
\usepackage{booktabs,tabularx,array,enumitem}
\usepackage{microtype}
\usepackage{xcolor}
\usepackage{fancyhdr}
\usepackage{hyperref}
\usepackage{tikz}
\usetikzlibrary{arrows.meta,positioning}
\definecolor{ink}{RGB}{29,53,80}
\definecolor{muted}{RGB}{92,100,110}
\hypersetup{colorlinks=true,linkcolor=ink,citecolor=ink,urlcolor=ink,
  pdftitle={Языки гипотез и сертификаты границ идентификации},
  pdfauthor={Александр Странд}}
\setlength{\parindent}{0pt}
\setlength{\parskip}{5pt plus 1.3pt minus .7pt}
\setlength{\emergencystretch}{3em}
\setlist[enumerate]{leftmargin=1.7em,topsep=3pt,itemsep=2pt}
\setlist[itemize]{leftmargin=1.7em,topsep=3pt,itemsep=2pt}
\pagestyle{fancy}\fancyhf{}
\fancyhead[L]{\footnotesize\color{muted}Языки гипотез и сертификаты границ идентификации}
\fancyhead[R]{\footnotesize\color{muted}А. Странд}
\fancyfoot[C]{\footnotesize\color{muted}\thepage}
\renewcommand{\headrulewidth}{.3pt}
\newtheorem{definition}{Определение}[section]
\newtheorem{theorem}{Теорема}[section]
\newtheorem{proposition}{Утверждение}[section]
\newtheorem{remark}{Замечание}[section]
\newcommand{\F}{\mathbb{F}}
\newcommand{\sem}[1]{\llbracket #1\rrbracket}
\newcommand{\VS}{\widetilde V}
\newcommand{\Agree}{\operatorname{Agree}}
\newcommand{\Bs}{\operatorname{Bs}}
\newcommand{\DFA}{\operatorname{DFA}}
\newcommand{\oneDCA}{\operatorname{1\text{-}DCA}}
\newcommand{\vstatus}[1]{\textsf{#1}}
\title{\bfseries\color{ink}Языки гипотез и сертификаты границ идентификации\\[4pt]
\large\normalfont Семантические пространства версий, локальные наблюдатели и конечные свидетельства невозможности ограниченных автоматов}
\author{Александр Странд}
\date{23 ноября 2026 г.}

\begin{document}
\maketitle
\vspace{-1.0em}
\begin{abstract}
\noindent Исследуется, какие выводы об идентификации гипотез допустимы при фиксированном языке, его расширении и изменении доступных наблюдений. Семантическое пространство версий отделено от синтаксических реализаций, а для вложенных линейных пространств булевых функций дано точное окно идентификации. Множество единогласия различает запросы, разделяющие кандидатов, и запросы, способные опровергнуть весь текущий класс. Рассмотрены два вычислительных свидетельства: конечный расширенный клеточный автомат с основным слоем Rule~110 и ограниченный синтез булевых языков с переходом между степенями. Основное дополнение настоящей редакции~--- \emph{конечный сертификат исключения}: проверенная матрица Майхилла--Нероуда размера $(k+1)\times(k+1)$ исключает \emph{все} DFA с не более чем $k$ состояниями как реализации полной функции на бесконечном домене. Этот вывод не означает нерегулярности произвольного чёрного ящика; для конкретно заданного языка $\{0^n1^n:n\ge0\}$ отдельно имеется Lean-доказательство отсутствия DFA любого конечного размера и корректности явного односчётчикового автомата. В работе явно отделяются конечные вычисления, доказанные свойства формальных моделей и утверждения, не перенесённые на исполняемый код.
\end{abstract}

\section{Семантическая идентификация и смена языка}

Для анализа обоснования будем использовать запись $\mathcal I=\langle\Phi,\mathcal D,\mathcal T,\mathcal E\rangle$: утверждение, его область и предпосылки, класс преобразований и тип проверяемого свидетельства. Само наличие такой записи не повышает доказательный статус.

\begin{definition}[Семантическое пространство версий]
Пусть $Q$~--- множество допустимых запросов, $A$~--- пространство ответов, $\mathcal L$~--- язык программ, $\sem\cdot:\mathcal L\to A^Q$~--- интерпретация, а $D\subseteq Q\times A$~--- конечный набор наблюдений. Тогда
\begin{equation}\label{eq:version}
\VS(\mathcal L,D)=\{\sem h:h\in\mathcal L,\ \forall(q,a)\in D\ \sem h(q)=a\}\subseteq A^Q.
\end{equation}
Программы с одинаковым поведением на всём заявленном $Q$ задают одну семантическую гипотезу, даже если их исходные тексты различаются.
\end{definition}

При фиксированной семантике новые непротиворечивые наблюдения могут только уменьшать пространство версий. Для согласованного вложения языков~--- только расширять его:
\begin{equation}\label{eq:monotonic}
\begin{aligned}
D\subseteq D'&\Longrightarrow\VS(\mathcal L,D')\subseteq\VS(\mathcal L,D),\\
\mathcal L_i\hookrightarrow\mathcal L_j&\Longrightarrow\VS(\mathcal L_i,D)\subseteq\VS(\mathcal L_j,D).
\end{aligned}
\end{equation}
Следовательно, единственность $|\VS(\mathcal L_i,D)|=1$ не обязана сохраняться в более широком $\mathcal L_j$. Статус \vstatus{IDENTIFIED} всегда включает ограничение на язык и семантический домен.

Две разные задачи идентификации выражаются перестановкой кванторов:
\begin{align}\label{eq:quantifiers}
&\forall i\ \forall f^\star\in\sem{\mathcal L_i}\ \exists D_i\text{ конечное}:\quad
 \VS(\mathcal L_i,D_i)=\{f^\star\};\\
&\exists D\text{ конечное}\ \forall i\ (f^\star\in\sem{\mathcal L_i})\Rightarrow
 \VS(\mathcal L_i,D)=\{f^\star\}.
\end{align}
Ни одно из этих утверждений \emph{не постулируется} для произвольной бесконечной иерархии. Первое может выполняться для каждого конечного ограниченного языка, но не влечёт второе с единым $D$.

Пример: на $Q=\F_2^2$ ответы $f(0,0)=0$, $f(1,0)=1$ выделяют $f=x_1$ среди функций, зависящих только от $x_1$, но оставляют \emph{четыре} полные булевы функции при расширении до всех $f:Q\to\F_2$. Количество AST и количество различных поведений на $Q$ при этом не совпадают.

\section{Двустороннее окно в линейных булевых языках}

Пусть $Q=\F_2^n$ и $R=\F_2[x_1,\ldots,x_n]/(x_1^2-x_1,\ldots,x_n^2-x_n)$~--- алгебра всех булевых функций, представленных мультилинейными полиномами. Рассмотрим вложенные конечномерные линейные пространства $U_0\subseteq\dots\subseteq U_m\subseteq R$, носитель наблюдений $Z\subseteq Q$ и согласованные значения $y\in\F_2^Z$.

Обозначим оператор ограничения $\rho_Z:U_i\to\F_2^Z$ и ядро $K_i=\ker(\rho_Z|_{U_i})=U_i\cap I(Z)$, где $I(Z)$~--- идеал функций, обращающихся в нуль на $Z$.

\begin{theorem}[Окно идентификации]\label{thm:window}
Если уравнение $\rho_Z(f)=y$ разрешимо в $U_i$, его пространство решений есть аффинный класс $f_i+K_i$, причём
\begin{equation}\label{eq:affine}
|\VS(U_i,D)|=2^{\dim U_i-\operatorname{rank}(\rho_Z|_{U_i})}.
\end{equation}
Пусть $i_*(D)$~--- первый индекс разрешимости, а $k_{\rm break}(Z)$~--- первый индекс с $K_i\ne\{0\}$; для отсутствующего индекса принимаем $+\infty$. Тогда
\begin{equation}\label{eq:window}
|\VS(U_i,D)|=1\quad\Longleftrightarrow\quad i_*(D)\le i<k_{\rm break}(Z).
\end{equation}
\end{theorem}
\begin{proof}
При разрешимости прообраз $y$ под линейным отображением $\rho_Z$ является сдвигом его ядра. Вложенность пространств делает разрешимость монотонной по $i$, а ядра~--- вложенными. Единственный элемент существует ровно после появления решения и до первого ненулевого ядра.
\end{proof}

Если аффинное пространство версий непусто, то запрос $q\in Q\setminus Z$ делит его поровну между ответами $0$ и $1$ \emph{тогда и только тогда}, когда $p\mapsto p(q)$ является ненулевым линейным функционалом на $K_i$. Это свойство линейной модели нельзя автоматически распространять на произвольные грамматики.

\section{Различение гипотез и проверка модели}

\begin{definition}[Множество единогласия]
Для \emph{непустого} $V\subseteq A^Q$ положим
\begin{equation}\label{eq:agree}
\Agree(V)=\{q\in Q:\forall f,g\in V\ \ f(q)=g(q)\}.
\end{equation}
\end{definition}

При $q\notin\Agree(V)$ гипотезы предлагают разные исходы: измерение способно разделить кандидатов. При $q\in\Agree(V)$ все кандидаты предсказывают один ответ; противоположный ответ, если измерение достоверно и допустимо, даёт $\VS(\mathcal L,D\cup\{(q,a)\})=\varnothing$ и отвергает \emph{текущий язык вместе с его предпосылками}. В общем случае при $|A|>2$ даже различающий запрос может вернуть исход, не предсказанный ни одной гипотезой.

\begin{figure}[htbp]
\centering
\begin{tikzpicture}[>=Latex,font=\footnotesize,node distance=6mm]
\node[draw=ink,rounded corners,align=center,text width=4.4cm,minimum height=1.15cm] (v) {Непустое пространство версий\\$V=\VS(\mathcal L,D)$};
\node[draw=ink!65,rounded corners,align=center,text width=4.15cm,minimum height=1.30cm,right=1.2cm of v,yshift=9mm] (d) {$q\notin\Agree(V)$\\разделение кандидатов};
\node[draw=ink!65,rounded corners,align=center,text width=4.15cm,minimum height=1.30cm,right=1.2cm of v,yshift=-9mm] (f) {$q\in\Agree(V)$\\проверка общей модели};
\draw[->,thick,ink] (v.east) -- ++(5mm,0) |- (d.west);
\draw[->,thick,ink] (v.east) -- ++(5mm,0) |- (f.west);
\end{tikzpicture}
\caption{Один формальный интерфейс запросов допускает две разные цели проверки. Наличие единственной гипотезы не заменяет проверки совместимости языка с измерениями.}\label{fig:query}
\end{figure}

При $V=\varnothing$ кванторы в~\eqref{eq:agree} истинны вакуозно; интерпретировать такую ситуацию как успешное единогласие запрещено. Кроме того, несогласованность может означать неадекватность языка, повреждение измерений или ошибку привязки контекста. Сам факт $|V|=0$ не идентифицирует причину.

\subsection{Расширение языка как проверяемое действие}

Для преобразований $T_i:\mathcal L\mapsto\mathcal L_i$ можно определить множества допустимых ответов при наблюдениях $D$:
\begin{equation}\label{eq:cross}
P_i(q\mid D)=\{f(q):f\in\VS(\mathcal L_i,D)\}.
\end{equation}
Если $P_i(q\mid D)$ и $P_j(q\mid D)$ непусты и не пересекаются, достоверный ответ $a$ на $q$ устраняет по меньшей мере один из двух конкурирующих языков. Это проверка \emph{способов расширения представления}, а не только отдельных программ. Порядок двух преобразований может менять язык: $T_i(T_j(\mathcal L))$ не обязан совпадать с $T_j(T_i(\mathcal L))$; конкретная некоммутативность требует предъявленного свидетельства, а не следует из одних названий операторов.

Минимальный префикс несогласованности задаётся, если он существует, формулой
\begin{equation}\label{eq:failprefix}
t_{\rm fail}(\mathcal L)=\min\{t:\VS(\mathcal L,D_{\le t})=\varnothing\}.
\end{equation}
Его сертификат требует проверки непустоты на предыдущем префиксе и отсутствия кандидатов на текущем; конечный перебор или корректное UNSAT-доказательство должны быть привязаны к версии грамматики. Более сложный язык не следует считать истинным только потому, что узкий опровергнут.

\section{Конечная локальная конструкция с основным слоем Rule 110}

Клеточный эксперимент использует периодическое кольцо из $64$ клеток. Основной бинарный слой до публикации ответа следует Rule~110; роли клеток, фазы, сообщения и регистры образуют \emph{расширенный} локальный алфавит. Поэтому речь не идёт о реализации всего наблюдателя исключительно двухсостояниевым Rule~110.

На четырёх специально заданных историях с двумя контекстами основные конфигурации каждой пары совпадают целиком на тактах $10\ldots27$, а дополнительная память сохраняет различие. В состоянии $t=13$ доступна контрольная чётность, с $t=19$~--- квитанция наблюдений; на переходе $27\to28$ исторический бит записывается обратно в основной слой. Полные состояния \emph{расширенного} автомата не сливаются, поэтому детерминизм не нарушается.

Допустимые тройки $C_3=\{000,010,110,111\}$ дополнены паритетным битом $p=z_1\oplus z_A\oplus z_B$:
\begin{equation}\label{eq:code}
C_4=\{0000,0101,1100,1111\},\qquad d_{\min}(C_4)=2.
\end{equation}
Такой код обнаруживает любую инверсию \emph{одного бита} корректно сформированного слова, но не гарантирует обнаружение двухбитной согласованной порчи: $10$ из $24$ двухбитных изменений ведут к другому допустимому слову. Нулевой синдром не удостоверяет происхождение произвольной начальной 64-битной конфигурации.

\begin{equation}\label{eq:live}
\begin{aligned}
\mathrm{Cert}_{\rm live}(V,a,z_1,z_A,z_B,p)\ \Longleftrightarrow\quad
 &[V=2^a]\\[-2pt]
 &\land[(z_1,z_A,z_B)\in C_3]\\[-2pt]
 &\land[h_{\rm code}(z_1,z_A,z_B)=a]\\[-2pt]
 &\land[z_1\oplus z_A\oplus z_B\oplus p=0].
\end{aligned}
\end{equation}

В отличие от однажды установленного флага \cite{leucker2009}, этот предикат относится к \emph{текущим} значениям памяти. Он не аутентифицирует произвольно согласованную подмену сразу всех полей. Статус $\vstatus{ORIGIN\_MEMBERSHIP}$ для произвольной исходной конфигурации остаётся $\vstatus{UNKNOWN}$.

\begin{table}[htbp]
\centering\small
\caption{Область формальных и вычислительных результатов для конечной клеточной конструкции}\label{tab:lean}
\renewcommand{\arraystretch}{1.10}
\begin{tabularx}{\linewidth}{@{}>{\raggedright\arraybackslash}p{3.2cm}>{\raggedright\arraybackslash}X>{\raggedright\arraybackslash}p{2.6cm}@{}}
\toprule
\textbf{Группа} & \textbf{Объект проверки} & \textbf{Основание}\\
\midrule
Локальные свойства & Алгебраическое правило 110, код $C_4$, конечный предикат $\mathrm{Cert}_{\rm live}$ & 10 теорем Lean\\
64-клеточный мост & Основной слой, слияние четырёх историй, часы и шесть регистровых вмешательств & 5 теорем Lean\\
Наблюдавшиеся состояния & Формальный размер алфавита, конечные траектории и некоторые классы повреждений & Протокол эксперимента\\
\bottomrule
\end{tabularx}
\end{table}

Предоставленные журналы Lean~4.36.0 фиксируют успешное завершение обоих модулей без \texttt{sorry}. Для пяти пространственно-временных теорем указан стандартный \texttt{propext}; доказательства относятся к перечисленным \emph{четырём} историям и заявленным конечным моментам. Они не устанавливают корректность для всех $2^{64}$ начальных состояний или для произвольного времени. В оригинальной редакции приведены $23\,014\,938\,240$ формальных комбинаций полей клетки и $505$ наблюдавшихся состояний на номинальных траекториях. Первое число~--- оценка декартова произведения полей, второе~--- статистика конкретного прогона, а не минимальность автомата.

Условная нижняя граница памяти тоже зависит от области: если один и тот же доступный контекст должен различать \emph{пять попарно различных} результатов, нужны не менее пяти различимых внутренних состояний. Из слияния \emph{двух} историй при фиксированном контексте такая пятисостоянийная граница для всей конструкции самостоятельно не следует.

\section{Конечная эскалация булева языка}

В отдельном Python-прототипе \cite{jha2017} \texttt{gate6\_autonomous\_escalation.py} рассматривается $Q=\F_2^3$ (восемь входов) и заранее объявленная иерархия языков алгебраических нормальных форм степеней $\le1,\le2,\le3$. Контроллер выбирает различающие запросы, проверяет оставшуюся единственную гипотезу на ещё не наблюдённых точках и, если $|\VS(\mathcal L_i,D)|=0$, переходит к следующему уровню.

Пять встроенных тестов прототипа выполнены успешно: цели $x_1x_2$, $x_1x_2x_3$ и $x_1\oplus x_3$ разделяют соответственно одну, две и ноль эскалаций; блокировка расширения предотвращает ложный сертификат, а изменение сохранённой чётности или добавление ложной записи об опровержении приводит к отказу локального проверяющего. Это \emph{конечная} демонстрация выбора степени при заранее заданной грамматике, а не изобретение неограниченного нового языка.

Проверяющий пересчитывает пространство версий по текущим наблюдениям и проверяет отсутствие моделей для \emph{имеющихся} записей об опровержении. Однако список таких записей реализован изменяемым Python-списком и не является криптографически неизменяемым журналом; XOR по записанным ответам сам по себе не является независимым удостоверением истинности измерений. Полное покрытие восьми входов перед выдачей результата не доказывает перенос на бесконечный домен.

\section{Gate 11: конечное свидетельство глобальной невозможности ограниченных DFA}

Здесь появляется более сильный тип отрицательного сертификата. Его область~--- \emph{бесконечное} множество слов $Q=\Sigma^*$, но исключаемый класс имеет \emph{фиксированный конечный ресурсный предел}.

Пусть $\Sigma=\{0,1\}$, $k\ge1$, а $f:\Sigma^*\to\{0,1\}$~--- целевая функция принадлежности (или достоверный оракул). По критерию Майхилла--Нероуда \cite{myhill1957,nerode1958} выберем префиксы $p_i=0^i$ и суффиксы $s_j=1^j$ для $0\le i,j\le k$. Их таблица ответов образует конечную матрицу
\begin{equation}\label{eq:mn_matrix}
M_k(i,j)=f(p_is_j)=f(0^i1^j),\qquad 0\le i,j\le k.
\end{equation}
\begin{theorem}[Конечный сертификат исключения DFA с не более чем $k$ состояниями]\label{thm:bounded_dfa}
Если все $(k+1)^2$ проверенных ответов удовлетворяют
\begin{equation}\label{eq:mn_diagonal}
M_k(i,j)=\begin{cases}1,&i=j,\\0,&i\ne j,\end{cases}
\end{equation}
то \emph{ни один} DFA с числом состояний $\le k$ не вычисляет $f$ на всей $\Sigma^*$.
\end{theorem}
\begin{proof}
Предположим, что такой DFA существует. После чтения $k+1$ префиксов $p_0,\ldots,p_k$ два из них, $p_i$ и $p_j$ при $i\ne j$, обязаны приводить автомат к одному состоянию. Поэтому добавление одного и того же суффикса $s_i=1^i$ даёт равные ответы. Однако по~\eqref{eq:mn_diagonal} они равны $f(p_is_i)=1$ и $f(p_js_i)=0$. Противоречие.
\end{proof}

Этот сертификат содержит конечные данные, но делает вывод \emph{о каждом слове потенциальной реализации DFA}, поскольку равенство двух его состояний имеет универсальное следствие для всех продолжений. Таким образом, отрицательный результат не ограничен длиной запрошенных слов. Предпосылка принципиальна: матрица должна соответствовать \emph{истинным} ответам того же целевого $f$ при неизменной привязке к контексту.

\begin{figure}[htbp]
\centering
\begin{tabular}{c|cccc}
$M_3$ & $s_0$ & $s_1$ & $s_2$ & $s_3$\\ \hline
$p_0$ & $\mathbf 1$ & 0 & 0 & 0\\
$p_1$ & 0 & $\mathbf 1$ & 0 & 0\\
$p_2$ & 0 & 0 & $\mathbf 1$ & 0\\
$p_3$ & 0 & 0 & 0 & $\mathbf 1$
\end{tabular}
\caption{Иллюстрация ограниченного сертификата при $k=3$: шестнадцать ответов исключают все DFA с не более чем тремя состояниями, но не DFA большего размера.}\label{fig:matrix}
\end{figure}

\begin{remark}[Принципиальное ограничение конечного оракула]\label{rem:oracle}
Для каждого \emph{фиксированного} $k$ матрица~\eqref{eq:mn_diagonal} исключает только $\DFA(\le k)$. Из любого конечного набора запросов без дополнительных предпосылок \emph{не следует} отсутствие DFA произвольного конечного размера. Конечная таблица ответов допускает продолжение до регулярного языка, совпадающего со всеми наблюдениями. Следовательно, такой эксперимент не доказывает нерегулярность неизвестного чёрного ящика.
\end{remark}

Для \emph{формально заданного} языка
\begin{equation}\label{eq:anbn}
L_{=}=\{0^n1^n:n\ge0\}
\end{equation}
имеется другой тип результата. В открытом модуле \href{https://github.com/Dnarts/proof-carrying-discovery/blob/main/formal/DFAto1DCASeparation.lean}{\texttt{DFAto1DCASeparation.lean}} приведены теоремы \texttt{anbn\_not\_dfa} для \emph{каждого} $S\in\mathbb N$ и \texttt{anbn1DCA\_accepts} для явно заданного детерминированного автомата с одним неограниченным счётчиком и тремя состояниями управления. В журнале Lean~4.36.0 указан код завершения $0$ и стандартные аксиомы \texttt{propext}, \texttt{Classical.choice}, \texttt{Quot.sound} (в зависимости от теоремы). Здесь квантор $\forall S$ обеспечен \emph{математическим доказательством о точно определённой функции}, а не набором конечных ответов внешнего оракула.

Отдельный Python-автомат, подобранный по конечным данным, не становится формально проверенным распознавателем всех слов лишь благодаря существованию этой Lean-теоремы. Для переноса необходима проверенная семантическая связь реализации с формальным определением.

\section{Типы свидетельств и доказательный статус}

\begin{table}[htbp]
\centering\small
\caption{Различные объекты сертификации на бесконечном домене}\label{tab:status}
\renewcommand{\arraystretch}{1.12}
\begin{tabularx}{\linewidth}{@{}>{\raggedright\arraybackslash}p{2.55cm}>{\raggedright\arraybackslash}X>{\raggedright\arraybackslash}p{3.05cm}@{}}
\toprule
\textbf{Статус} & \textbf{Что удостоверяется} & \textbf{Не следует}\\
\midrule
Исключён класс & Нет DFA с $\le k$ состояниями, согласованного со всем целевым $f$; матрица достоверна & Нерегулярность оракула\\
Согласовано с $D$ & Реализация объясняет конечную выборку & Полная правильность\\
ID внутри $G$ & Один семантический кандидат относительно фиксированного $G$ & Уникальность после смены $G$\\
Формальная сепарация & Для заданного $L_=$ доказаны $L_=\notin\DFA$ и $L_=\in\oneDCA$ & Корректность любого Python-кода\\
\bottomrule
\end{tabularx}
\end{table}

Семантическая идентичность не равна равенству исходных программ. Эквивалентность наблюдательных поведений не равна идентичности физических причин. Историческая квитанция не гарантирует сохранности текущего регистра; проверка синдрома не заменяет доказательства происхождения данных. Эти границы могут быть включены в проверяемый протокол, но не устраняются присвоением нового названия статусу.

Отдельно важна асимметрия Gate~11: \emph{опровержение всего ресурсно ограниченного класса} может быть подтверждено конечным сертификатом, тогда как \emph{принятие новой реализации по конечной выборке} остаётся относительным. Оба принципа опираются на известную теорию конечных автоматов и идентификации; научное содержание данного примера~--- в совместном предъявлении формулировки, области, проверяемого свидетеля и ограничений переноса.

\section{Заключение}

Семантическая единственность зависит от класса допустимых гипотез. Запросы различают модели либо обнаруживают их коллективную несовместность с измерением. Свидетельство невозможности для всех DFA с числом состояний $\le k$ может быть \emph{конечным} и одновременно иметь \emph{универсальную силу внутри указанного ограниченного класса}. Полное доказательство нерегулярности конкретного языка требует отдельного квантора по всем $k$ и точной семантики целевого объекта. Разделение этих утверждений позволяет сохранять отрицательные результаты при смене представления без необоснованного повышения статуса положительных результатов.

Клеточно-автоматная конструкция и конечный прототип эскалации булевой степени иллюстрируют операционное применение этих ограничений, но сами по себе не доказывают открытую способность произвольной вычислительной системы порождать новые научные языки.

\begin{thebibliography}{99}\setlength{\itemsep}{1pt}
\bibitem{myhill1957} J.~Myhill. \emph{Finite Automata and the Representation of Events}. WADC Technical Report 57--624, 1957.
\bibitem{nerode1958} A.~Nerode. \emph{Linear Automaton Transformations}. Proceedings of the American Mathematical Society, 9(4):541--544, 1958.
\bibitem{moore1956} E.~F.~Moore. \emph{Gedanken-experiments on Sequential Machines}. Automata Studies, pp.~129--153, 1956.
\bibitem{gold1978} E.~M.~Gold. \emph{Complexity of Automaton Identification from Given Data}. Information and Control, 37(3):302--320, 1978.
\bibitem{jha2017} S.~Jha, S.~A.~Seshia. \emph{A Theory of Formal Synthesis via Inductive Learning}. Acta Informatica, 54:693--726, 2017.
\bibitem{hamming1950} R.~W.~Hamming. \emph{Error Detecting and Error Correcting Codes}. Bell System Technical Journal, 29(2):147--160, 1950.
\bibitem{leucker2009} M.~Leucker, C.~Schallhart. \emph{A Brief Account of Runtime Verification}. Journal of Logic and Algebraic Programming, 78(5):293--303, 2009.
\end{thebibliography}
\end{document}
